
\thispagestyle{empty}
\begin{center}
    Wykaz obszarów i narzędzi GenAI \\
    wykorzystanych przez autora w trakcie wykonywania pracy dyplomowej
\end{center}

\begin{table}[ht]
    \centering
    \begin{tabularx}{\linewidth}{|c|X|X|}\hline
       \textbf{Lp.}  & \multicolumn{1}{c|}{\textbf{Obszar wykorzystywania $^{a)}$}} & \multicolumn{1}{c|}{\textbf{Narzędzie}} \\ \hline
         &  &\\ \hline
         &  &\\ \hline
    \end{tabularx}
    \label{tab:my_label}
\end{table}

\textbf{Uwaga!} W przypadku, gdy autor nie korzysta z obszarów i narzędzi GenAI w tabeli wpisuje „nie korzystano”

\

\textit{$^{a)}$ Przykładowe obszary wykorzystania:
    \begin{enumerate}[label=\alph*)]
        \item redakcja i korekta tekstu (m.in. poprawa stylistyczna i gramatyczna, synonimy i~parafrazowanie, sprawdzenie spójności językowej),
        \item analiza stanu wiedzy, przegląd literatury,
        \item generowanie treści lub przykładów,
        \item tworzenie pytań do ankiet i kwestionariuszy,
        \item tworzenie schematów, diagramów i map myśli,
        \item analiza danych i wizualizacja wyników (m.in. obliczenia statystyczne, tworzenie wykresów i tabel),
        \item tworzenie podsumowań i wniosków,
        \item pisanie i testowanie kodu,
        \item tłumaczenia językowe.
    \end{enumerate}}